\documentstyle[11pt,newsymb]{article}
\begin{document}

\title{A Short Proof of a Conjecture on the Integral Means of the
Derivative of a Convex Function}

\author{Roger W.\ Barnard\thanks{Research supported in part by
Texas Advanced Research Program grant \#003644-125.} and Kent
Pearce\\
Department of Mathematics, Texas Tech University}

\date{}

\maketitle

For $d>0$ let $D_d=\{z:|z|<d\}$ with $D_1=D$ and let $\partial D_d$ denote
the boundary of $D_d$.  Let $S$ be the standard class of analytic,
univalent functions $f$ on $D$, normalized by $f(0)=0$ and $f'(0)=1$
and let $K$ denote the well-known class of convex functions in $S$.
For $0\leq\alpha<1$ let $S^\ast(\alpha)$ denote the subclass of $S$
of starlike functions of order $\alpha$, i.e., a function $f\in
S^\ast(\alpha)$ if and only if $f$ satisfies the condition ${Re}\:
zf'(z)/f(z)>\alpha$ on $D$.  It is well known that $K\subset S^\ast
(1/2)$.

For $F\subset S$ and for $1/4\leq d\leq 1$ let
$$F_d=\{f\in F:\min_{z\in D}\left|\frac{f(z)}{z}\right|=d\}.$$
Note, that $K_d=\emptyset$ for $1/4\leq d<1/2$.

A general problem which arose out the authors' work in the early 80's
on omitted value problems for convex functions, see \cite{Ba} and
\cite{BaPe1}, is the following: Given $F\subset S$ and $1/4\leq d\leq
1$ determine the sharp constant $A=A(F_d)$ such that for any $f\in F_d$
\begin{equation}
\label{eq:1}
I_{-1}(f')=\frac{1}{2\pi}\int^{2\pi}_0
\left|\frac{1}{f'(e^{i\theta})}\right|d\theta\leq\frac{A}{d}.
\end{equation}
Further, determine the sharp constant $A=A(F)=\displaystyle{\sup_d}
A(F_d)$.

It follows fairly easily from subordination theory that
$A(S^\ast(1/2))\leq 4/\pi$.  Furthermore, this estimate is sharp for
$S^\ast(1/2)$ since the functions $f_n(z)=z/(1-z^n)^{1/n}$ belong to
$S^\ast(1/2)$ for each $n>0$.  However, this estimate is not sharp
for the class $K$ of convex functions which is a proper subset of
$S^\ast(1/2)$.  Considerable numerical evidence suggested to the
authors to make the following conjecture.

{\bf Conjecture.} For each $d,\: 1/2\leq d\leq 1,\: A=A(K_d)=1$ in
(1) with equality holding for all domains which are bounded by
regular polygons centered at the origin.

This conjecture was announced in March 1985 at the Symposium on the
Occasion of the Proof of the Bieberbach Conjecture at Purdue
University.  It also appeared as Conjecture 8 in the first author's
``Open Problems and Conjectures in Complex Analysis'' in \cite{Ba}.  It
was thought, by many function theorists, that the conjecture would be
easily settled, given the vast literature on convex functions and the
large research base for determining integral mean estimates, see
\cite{Du}.

An initial difficulty was the non-applicability of Baerenstein's
circular symmetrization methods, since convexity, unlike univalence
and starlikeness, is not preserved under circular symmetrization.
Although Steiner symmetrization does preserve convexity, see \cite{Va}, it
did not appear to be helpful for the problem and, indeed, we found
that extremal domains need possess no standard symmetry.

A confusing issue, which also arises, is that the integral means of
the standard approximating functions $f_n$ in $K$ defined by
$$f'_n(z)=\prod^n_{k=1}(1-ze^{i\theta_k})^{-2\alpha_k},0<\alpha_k\leq
1,\sum^n_{k=1}\alpha_k=1$$ 
decrease when the arbitrarily distributed $\theta_k$ are replaced by
uniformly distributed $t_k=k\pi/n$, as was shown in \cite{Zh}.  The
conjecture suggests that multiplication by the minimum modulus $d$
must overcome this decrease.

We make the following definition.

{\bf Definition.}  Let $\Gamma$ be a curve in $\Bbb C$ such that the
left- and right-hand tangents to the curve $\Gamma$ exist at each
point on $\Gamma$.  The curve $\Gamma$ will be said to {\em
circumscribe a circle} $C$ if the left- and right-hand tangents to
the curve $\Gamma$ at each point on $\Gamma$ lie on tangent lines to
the circle $C$.

We will employ the following notation.

{\bf Notation.} Let $f\in S$ and suppose that $\gamma$ is a subarc of
$\partial D$ on which $f$ is smooth.  For $z=e^{i\theta}\in\gamma$
let $d_{\theta}=<f(z)$, $\frac{zf'(z)}{|zf'(z)|}>$, i.e, $d_\theta$
is the directed length of the projection of $f(z)$ onto the outward
unit normal to $\partial f(D)$ at $f(z)$.

In 1993 \cite{BaPe2}, we proved the following theorem which verified the
conjecture.

{\bf Theorem A.} Let $f\in K$, $d=\displaystyle{\min_\theta}
|f(e^{i\theta})|$ and $d^\ast=\displaystyle{\sup_\theta}\; d_\theta$.
Then,
\begin{equation}
\label{eq:2}
\frac{1}{d^\ast}\leq\frac{1}{2\pi}\int^{2\pi}_{0}\frac{1}
{|f'(e^{i\theta})|}d\theta\leq\frac{1}{d}
\end{equation}
with equality holding if $\partial f(D)$ circumscribes $\partial
D_d$.

The original proof, which was lengthy, was based on the Julia
variational formula and depended heavily on the convexity of $f$.  We
obtained, arising out of the proof, the rather unexpected sufficient
condition for equality to occur in (2) for the classes $K_d$.
However, because the proof used a scheme to approximate convex
functions by polygonally convex functions, we did not obtain a
necessary condition for equality.

We have devised a new, simpler proof for the conjecture which extends
Theorem A.  The proof releases the convexity requirement and
validates the necessity of the sufficient condition.

{\bf Theorem B.} Let $f\in S^\ast(\alpha)$ for some $\alpha>0$.
Suppose $f$ is smooth on $X\subset\partial D$ where $X$ is a
countable union of pairwise disjoint subarcs of $\partial D$ such
that the complement of $X$ in $\partial D$ has measure zero.  Let
$d_\ast=\displaystyle{\inf_{\theta\in X}}\:\: d_\theta$,
$d^\ast=\displaystyle{\sup_{\theta\in X}}\:\: d_\theta$.  Then,
\begin{equation}
\label{eq:3}
\frac{1}{d^\ast}\leq\frac{1}{2\pi}\int^{2\pi}_0
\frac{1}{|f'(e^{i\theta})|}d\theta\leq\frac{1}{d_\ast}
\end{equation}
with equality holding if and only if $\partial f(D)$ circumscribes
$\partial D_{d_\ast}$.

{\bf Remark.} In general $d_\ast\leq d=\displaystyle{\min_{z\in
D}|\frac{f(z)}{z}|}$ with equality if $f\in K$.

{\bf Proof.} Let $X=\bigcup^\infty_{k=1}\gamma_k$, where each
$\gamma_k$ is a subarc of $\partial D$.  We have from the Cauchy
Integral Formula, with $z=e^{i\theta}$, that
\begin{eqnarray}
\nonumber
& 1 & =\lim_{r\rightarrow 1}\frac{1}{2\pi}\int^{2\pi}_0  
\frac{f(rz)}{rzf'(rz)}d\theta\\[2ex]
\nonumber
&& =\frac{1}{2\pi}\int^{2\pi}_0{Re}\frac{f(z)}{zf'(z)}
d\theta\\[2ex]
\label{eq:4}
&& =\frac{1}{2\pi}\int^{2\pi}_0\frac{|f(z)|\cos
(\arg\frac{f(z)}{zf'(z)})}{|zf'(z)|}d\theta\\[2ex]
\nonumber
&& =\frac{1}{2\pi}\sum^\infty_{k=1}\int_{\gamma_k}
\frac{<f(z),\:\frac{zf'(z)}{|zf'(z)|}>}{|f'(z)|}d\theta\\[2ex]
\nonumber
&& =\frac{1}{2\pi}\sum^\infty_{k=1}\int_{\gamma_k}
\frac{d_\theta}{|f'(z)|}d\theta
\end{eqnarray}
The restriction of $f$ to $S^\ast(\alpha)$ with $\alpha>0$ assures
the passage of the limit in (4).  Replacing $d_\theta$ by $d^\ast$ and
$d_\ast$ in this last integral gives the left- and right-hand
inequalities in (3), respectively.  Equality clearly occurs if
$\partial f(D)$ circumscribes $\partial D_{d_\ast}$, for in this case
$d_\ast=d_\theta=d^\ast$ on $X$.  Conversely, if $\partial f(D)$ does not
circumscribe $\partial D_{d_\ast}$, then there must exist a subarc $I$,
which must lie in one of the $\gamma_k$, of positive (linear) measure
on which 
\begin{equation}
\label{eq:5}
d_\ast<d_\theta<d^\ast.
\end{equation} 
The strictness of (5) on $I$ implies that strict inequality holds in
both the left- and right-hand inequalities in (3).  $\blacksquare$

The authors would like to thank Al Baerenstein for his helpful
suggestion which led to their shorten proof.


\begin{thebibliography}{9}

\bibitem[Ba]{Ba}  Barnard, R.W.\ ``Open problems and conjectures in
complex analysis'', in {\em Computational Methods and Function
Theory}, Lecture Notes in Math.\ No.\ 1435 (Springer-Verlag, 1990),
1--26.

\bibitem[BaPe1]{BaPe1}  Barnard, R.W.\ and Pearce, K.\ ``Rounding corners of
gearlike domains and the omitted area problem'', {\em J.\ Comput.\
Appl.\ Math.} {\bf 14} (1986), 217--226.

\bibitem[BaPe2]{BaPe2}  Barnard, R.W.\ and Pearce, K.\ ``Sharp bounds on the
$H_p$ means of the derivative of a convex function for $p=-1$'', {\em
Complex Analysis} {\bf 21} (1993), 149--158.

\bibitem[Du]{Du}  Duren, P.\ {\em Univalent Functions}, Springer-Verlag,
New York, 1980.

\bibitem[Va]{Va}  Valentine, F.A.\ {\em Convex Sets}, McGraw-Hill, New
York, 1964.

\bibitem[Zh]{Zh}  Zheng, J.\ ``Some extremal problems involving $n$ points
on the unit circle'', Dissertation, Washington University (St.\
Louis), 1991.

\end{thebibliography}

 \end{document}
